%% ===========================================================================
%%  Report template -- Exoplanet Transit Detection Winter Project
%%  Copy this file into your own repo (e.g. report/report.tex), fill it in,
%%  and compile with pdflatex + bibtex (or upload to Overleaf).
%%  Target length: 4 pages (freshmen/sophomores) to 6 pages (juniors/seniors).
%% ===========================================================================
\documentclass[11pt,a4paper]{article}
\usepackage[T1]{fontenc}
\usepackage[utf8]{inputenc}
\usepackage{lmodern}
\usepackage{amsmath,amssymb}
\usepackage[margin=2.2cm]{geometry}
\usepackage{graphicx}
\usepackage{booktabs}
\usepackage[numbers,sort&compress]{natbib}
\usepackage[colorlinks=true,linkcolor=blue,urlcolor=blue,citecolor=blue]{hyperref}

\newcommand{\Rearth}{\ensuremath{R_\oplus}}
\newcommand{\todo}[1]{\textcolor{red}{[TODO: #1]}}

\title{\todo{A title that states your result}\\[2pt]
       \large e.g.\ ``Recovering Kepler-10\,b and measuring the detection efficiency of a simple transit pipeline''}
\author{\todo{Your Name}\\ \small Astronomy Club Winter Project \quad
        \href{https://github.com/YOU/REPO}{github.com/YOU/REPO}}
\date{\today}

\begin{document}
\maketitle

\begin{abstract}
\todo{At most 150 words. One sentence each: the question; the data; the method; the main
results WITH NUMBERS (planet radius $\pm$ error; completeness or $S_{50}$); the one conclusion.}
\end{abstract}

\section{Introduction}
\todo{The transit method in two or three sentences, with Eq.~\eqref{eq:depth}. Why measuring the
completeness of a search matters (occurrence rates). What you did in this project, in one paragraph.}
\begin{equation}
  \delta \simeq \left(R_p/R_\star\right)^2 \label{eq:depth}
\end{equation}

\section{Data}
\todo{Mission, target(s), quarters/sectors, cadence, flux type (PDCSAP), number of points,
time span. Stellar parameters with their source in a table like Table~\ref{tab:star}.}
\begin{table}[h]
\centering
\caption{Stellar and published planetary parameters.}
\label{tab:star}
\begin{tabular}{@{}lll@{}}
\toprule
Parameter & Value & Source\\
\midrule
$R_\star$ [$R_\odot$] & \todo{} & \todo{e.g.\ NASA Exoplanet Archive / TIC}\\
$M_\star$ [$M_\odot$] & \todo{} & \\
Published $P$ [d] & \todo{} & \\
Published $R_p$ [\Rearth] & \todo{} & \\
\bottomrule
\end{tabular}
\end{table}

\section{Methods}
\subsection{Detrending}
\todo{Method, window length and why; gap handling; sigma clipping (upper only, why).
Second pass with transits masked.}
\subsection{Transit search}
\todo{BLS: period range, duration grid, oversampling; detection statistic and threshold.}
\subsection{Transit fit}
\todo{Model (trapezoid / batman), free and fixed parameters, limb darkening coefficients and
source, supersampling, priors; how uncertainties were estimated.}
\subsection{Injection--recovery}
\todo{Grid in period and radius, injections per cell, how $P$, $R_p$, $t_0$ were drawn,
injection model, and your exact recovery criteria. State that the pipeline settings were frozen.}

\section{Results}
\todo{Periodogram figure; folded transit with model; table comparing $R_p$ estimates with
literature; completeness map; efficiency curve with $S_{50}$; your controlled experiment.}
% \begin{figure}[h]
%   \centering
%   \includegraphics[width=0.8\linewidth]{../figures/fold.png}
%   \caption{Folded light curve of ... with the best-fit model. Note ...}
%   \label{fig:fold}
% \end{figure}

\section{Discussion}
\todo{Does your radius agree with the literature (within how many sigma)? If not, why?
What limits your pipeline, according to the completeness map? What would you change next,
and what would you expect to gain?}

\section{Conclusions}
\todo{Three bullet points.}

\paragraph{Acknowledgements.} This paper includes data collected by the Kepler/TESS mission,
obtained from the MAST data archive at the Space Telescope Science Institute (STScI).
This work used Lightkurve~\citep{Lightkurve2018}, Astropy~\citep{Astropy2022}, NumPy, SciPy and Matplotlib.
\todo{Mention people who helped, and how you used AI tools, if at all.}

\bibliographystyle{plainnat}
\bibliography{../docs/references}   % or your own .bib file

\end{document}
